\documentclass[11pt]{article}
% === core_packages.tex ===
\usepackage{amsmath, amssymb, amsthm, amsfonts}
\usepackage{geometry}
\usepackage{hyperref}
\usepackage{graphicx}
\usepackage{physics}
\usepackage{booktabs}
\usepackage{listings}
\usepackage{bookmark}
\usepackage{amscd}
\usepackage{tikz-cd}
\usepackage{mathtools}
\usepackage{xspace}
\usepackage[backend=biber, style=numeric]{biblatex}
\usepackage{tcolorbox}
\tcbuselibrary{theorems}
\usepackage{parskip}
\usepackage{tikz}
\usepackage{etoolbox}
\usetikzlibrary{arrows.meta, shapes, positioning, calc, cd, shapes.geometric, backgrounds}
\usepackage[en-US]{datetime2}
\usepackage{array}


% --- Math & Ontology Macros ---
\newcommand{\entropy}{S}
\newcommand{\opponent}{opponent}
\newcommand{\aside}[1]{[\emph{\opponent}: ``#1'']}
\newcommand{\paper}[1]{[\emph{#1}: ``#1'']}
\newcommand{\iame}{Real\textsuperscript{e}}
\newcommand{\iameverything}{Real\textsuperscript{everything}}
\newcommand{\fabric}{voxel}
\newcommand{\Fabric}{Voxel}
\newcommand{\pdflink}[1]{}



% --- Global Layout Defaults ---
\geometry{letterpaper, margin=1in}
\hypersetup{
    colorlinks=true,
    linkcolor=black,
    citecolor=black,
    filecolor=magenta,
    urlcolor=blue
}
\hfuzz=5pt \vfuzz=5pt \hbadness=10000 \vbadness=10000

\author{\iame}


\usepackage{tocloft}

% --- Article Specific Theorems ---
\newtheorem{lemma}{Lemma}
\theoremstyle{definition}
\newtheorem{axiom}{Axiom}
\newtheorem{theorem}{Theorem}
\newtheorem{corollary}{Corollary}
\newtheorem{definition}{Definition}
\newtheorem{proposition}{Proposition}
\newtheorem{conjecture}{Conjecture}


\makeatletter
\let\remark\@undefined
\let\endremark\@undefined
\makeatother
\theoremstyle{remark}
\newtheorem*{remark}{Remark}

\newtcbtheorem[number within=section]{principle}{Principle}%
{colback=blue!5,colframe=blue!50!black,fonttitle=\bfseries}{principle}

% --- Article TOC Layout ---
\setlength{\cftbeforesecskip}{1.0em}
\renewcommand{\cftsecleader}{\cftdotfill{\cftdotsep}}
\renewcommand{\cftsubsecleader}{\cftdotfill{\cftdotsep}}


\addbibresource{../references.bib}

\title{A Zero-Parameter Theory of Everything:\\
Existence, Compression, and the Emergence of Physical Law}

\author{Juha Meskanen}
\date{2026}

\begin{document}
\maketitle

% ============================================================
\begin{abstract}
We present a complete, zero-parameter derivation of physical reality from
a single principle: \emph{observers exist, and their probability of
existence is determined by the compressibility of their description}. From
this principle alone --- with no metric tensor, no action principle, no
Hilbert space, no cosmological constant, and no free parameters --- we
derive the structure of spacetime, the origin of quantum mechanics, and
the unification of General Relativity with quantum field theory.

The central equation is a Boltzmann-like probability measure over
histories $\gamma$ compatible with an observer $O$:
\[
\mathbb{P}(\gamma \mid O)
= \frac{1}{Z_O}\exp\!\bigl(-\mathcal{C}_O[\gamma]\bigr),
\qquad \gamma \in \Gamma_O,
\]
where $\mathcal{C}_O[\gamma]$ is the total description cost of history
$\gamma$ given observer $O$, and $Z_O$ is the partition function over
observer-compatible histories. Physical law, geometry, and quantum
behaviour are not postulated: they are the large-deviation minimisers of
$\mathcal{C}_O$.

The cost functional decomposes into three regimes --- discrete
information (\textbf{D}), spectral/quantum description ($\boldsymbol{\psi}$),
and geometric description (\textbf{G}) --- which we call the \textbf{D-$\psi$-G
trinity}. This trinity is not assumed; it is forced by the boundary
structure of any self-describing observer. We further show that the
apparent free parameter $n$ (the bit count of the universe) is itself
emergent: it is the saddle point of the observer's own probability of
existence. The value $n^* \approx 184$ predicted by this saddle-point
condition is consistent with the independently derived value from
inflationary cosmology, black hole thermodynamics, and the e-fold count
of inflation. The theory has no free parameters.
\end{abstract}

\textbf{Keywords:} theory of everything, Solomonoff induction, spectral
complexity, quantum gravity, information theory, observer self-selection,
Kolmogorov complexity, Wheeler--DeWitt, De Sitter, Bekenstein--Hawking

% ============================================================
\section{Introduction}
\label{sec:intro}

General Relativity and quantum mechanics each require an input the other
cannot supply. General Relativity takes the geometry of spacetime as
given and derives how matter curves it. Quantum mechanics takes the
Hilbert space structure of states as given and derives how they evolve.
Neither theory explains why its own input exists.

We argue both inputs are the same thing: a single integer $n$, the total
bit count of a finite, static, timeless universe, and the rule that
structures dominate the measure in proportion to how compressibly they
can be described.

From this starting point --- with no metric tensor, no action principle,
no cosmological constant, no Hilbert space postulated --- the following
emerge as counting consequences: the aspect ratio of spacetime, the
Bekenstein--Hawking entropy of a solar-mass black hole, the e-fold count
of inflation, the Friedmann equation, the wave-like character of matter,
and the classical limit of general relativity.

The universe does not obey quantum mechanics and general relativity. It
is compressed information, and those two theories are what compression
looks like from the inside. The only first principle is the existence of
an observer. Everything else follows.

% ============================================================
\section{The Single Principle and Its Consequences}
\label{sec:principle}

\subsection{The Observer as the First Principle}

We take as the unique starting point:

\begin{quote}
\textbf{Axiom (Observer Existence).} There exists at least one observer.
\end{quote}

An observer is defined operationally, following Paper~I
\cite{meskanen2001}: a finite information structure whose state sequence
encodes causally efficacious subjective experience. Paper~I establishes
from four modest axioms that such a structure must be finite, substrate
independent, and that the universe containing it must be static and
timeless. We take those results as established and build on them here.

\subsection{Probability of Existence via Solomonoff Induction}

Given the observer axiom, the natural question is: \emph{which universe
does the observer find herself in?} By the substrate independence of
Paper~I, the observer is a finite information structure, and the universe
is the minimal information structure that contains her. The probability
of a universe (history) $\gamma$ given observer $O$ is therefore
determined by the description length of $\gamma$:

\begin{equation}
\mathbb{P}(\gamma \mid O)
= \frac{1}{Z_O}\exp\!\bigl(-\mathcal{C}_O[\gamma]\bigr),
\qquad \gamma \in \Gamma_O,
\label{eq:master}
\end{equation}

where $\Gamma_O$ is the set of all histories compatible with the
existence of $O$, $\mathcal{C}_O[\gamma]$ is the total description cost
of $\gamma$ given $O$, and $Z_O = \sum_{\gamma \in \Gamma_O}
\exp(-\mathcal{C}_O[\gamma])$ is the partition function. This is
Solomonoff induction \cite{solomonoff1964} in path-integral form: the
prior over universes is exponentially suppressed by description length.

Equation~\eqref{eq:master} has one apparent free parameter, $\lambda$,
which we have set to unity without loss of generality: rescaling
$\lambda$ is equivalent to rescaling $\mathcal{C}_O$, and only the
ordering of histories by cost is physically meaningful. The parameter
$\lambda$ is therefore not physical. It is a units choice.

\subsection{Physical Law as Large-Deviation Minimiser}

The observed regularities of physics --- conservation laws, geometry,
quantum behaviour --- are not postulated in this framework. They are the
histories $\gamma^*$ that minimise $\mathcal{C}_O[\gamma]$ subject to
observer compatibility:

\[
\gamma^* = \arg\min_{\gamma \in \Gamma_O} \mathcal{C}_O[\gamma].
\]

A universe that is chaotic, infinite, or structureless is not forbidden
--- it simply has an astronomically higher description cost and therefore
contributes negligibly to the measure. We do not assume finiteness; we
derive it, because finite universes compress better than infinite ones.
We do not assume order; we derive it, because ordered histories are
shorter to describe than disordered ones.

% ============================================================
\section{The D-$\psi$-G Trinity}
\label{sec:trinity}

\subsection{Why Three Regimes?}

The cost functional $\mathcal{C}_O[\gamma]$ is not a single expression
across all scales. Any observer requires well-defined boundaries: she
must be distinguishable from her environment. The cost of specifying
those boundaries depends on the representation in which they are
described, and different representations dominate at different scales.

We identify three natural regimes:

\begin{itemize}
  \item \textbf{D} (Discrete information): The raw bit level. Boundaries
        are subsets of the bitstring. Cheap to define but carries no
        structure.
  \item \textbf{$\psi$} (Spectral/quantum): Boundaries are subspaces of
        a Hilbert space. Specifying a localised subspace requires many
        high-frequency modes; the spectral cost is high.
  \item \textbf{G} (Geometric): Boundaries are surfaces in a metric
        space. A closed surface in three dimensions is fully specified
        by a compact, low-frequency description. The geometric cost of
        a boundary is minimised relative to all other representations.
\end{itemize}

\subsection{Why Geometry Dominates at Observer Scales}

Observers are boundary-having structures. The representation that
minimises the boundary description cost dominates the measure via
equation~\eqref{eq:master}. Geometric boundaries --- surfaces --- are
cheaper to describe than Hilbert-space boundaries for any structure
above the Planck scale.

This is why we find ourselves in a geometric universe rather than a
Hilbert space. We do not postulate spacetime. We derive that geometric
descriptions dominate the measure for observers like us, because
specifying ``where Alice ends and the universe begins'' costs
exponentially fewer bits in geometric language than in spectral language.

\subsection{The Decomposition of $\mathcal{C}_O$}

The total cost decomposes as:

\begin{equation}
\mathcal{C}_O[\gamma]
= \underbrace{C_D[n]}_{\text{bit budget}}
+ \underbrace{C_\psi[\psi_\gamma]}_{\text{spectral cost}}
+ \underbrace{C_G[g_\gamma]}_{\text{geometric cost}},
\label{eq:decomp}
\end{equation}

where $C_D[n] = \log_2 n$ is the cost of specifying the bit count,
$C_\psi[\psi_\gamma]$ is the spectral complexity of the wavefunction
(defined in Paper~VI \cite{meskanen2026vi}), and $C_G[g_\gamma]$ is the
geometric cost of the history's spacetime configuration. The three
regimes are not independent: as shown in Paper~VII
\cite{meskanen2026vii}, minimising $C_\psi$ in Wheeler--DeWitt
minisuperspace recovers the same classical geodesics as minimising the
Euclidean action $S_E$ via Wick rotation. The spectral and geometric
costs are two representations of the same minimisation.

\subsection{The Three Faces of One Equation}

The same cost functional $\mathcal{C}_O$ has three limiting forms
depending on scale:

\begin{center}
\begin{tabular}{lll}
\hline
\textbf{Scale} & \textbf{Dominant cost} & \textbf{Emergent theory} \\
\hline
Planck scale & $C_D[n]$ & Discrete, quantised spacetime \\
Quantum scale & $C_\psi[\psi_\gamma]$ & Quantum mechanics, Born rule \\
Observer scale & $C_G[g_\gamma]$ & General Relativity, Friedmann \\
\hline
\end{tabular}
\end{center}

\noindent Quantum mechanics is not added to the framework: it is the
regime in which spectral cost dominates. General Relativity is not
added: it is the regime in which geometric boundary cost dominates.
Quantum gravity is not a separate theory to be discovered: it is
equation~\eqref{eq:master}, read at the scale where $C_\psi$ and $C_G$
are comparable.

% ============================================================
\section{The Emergence of $n$}
\label{sec:n}

\subsection{$n$ Is Not a Free Parameter}

The bit count $n$ appears in equation~\eqref{eq:decomp} as the argument
of $C_D[n] = \log_2 n$. One might ask whether $n$ is a free parameter
of the theory --- the one number that must be input by hand. We argue it
is not.

The probability that observer $O$ exists at all, marginalising over all
histories, is:

\begin{equation}
\mathbb{P}(O \mid n)
= \frac{Z_O(n)}{Z(n)},
\label{eq:pobs}
\end{equation}

where $Z(n) = \sum_{\gamma} \exp(-\mathcal{C}[\gamma])$ is the full
partition function over all histories of an $n$-bit universe. This
quantity depends on $n$: too small an $n$ and the universe cannot
contain a structure as complex as $O$; too large an $n$ and $O$ becomes
a negligible fluctuation within an overwhelmingly large configuration
space.

The prior over $n$ is itself determined by description length:
$\mathbb{P}(n) \propto 2^{-\log_2 n} = 1/n$. Larger values of $n$ are
penalised because they require more bits to specify. The joint
probability of observer existence at bit count $n$ is:

\begin{equation}
\mathbb{P}(O, n)
\propto \frac{Z_O(n)}{Z(n) \cdot n}.
\label{eq:joint}
\end{equation}

\subsection{The Saddle-Point Equation}

The observed $n$ is the value that maximises equation~\eqref{eq:joint}.
Taking the logarithmic derivative and setting it to zero:

\begin{equation}
\boxed{
\frac{\partial}{\partial n}
\Bigl[\log Z_O(n) - \log Z(n) - \log n\Bigr] = 0.
}
\label{eq:saddle}
\end{equation}

This is the self-consistency equation for the universe's bit count. It
has no free parameters. It asks: at what $n$ is the observer's
contribution to the partition function maximised relative to the total,
after penalising for the cost of specifying $n$ itself?

\subsection{Consistency with Independent Derivations}

Papers~IV and~V \cite{meskanen2026iv, meskanen2026v} derive $n$ by two
independent routes:

\begin{enumerate}
  \item \textbf{Inflationary calibration.} The aspect ratio
        $\mathcal{A}(n) = 2^n / (n \ln n)$ is matched to the
        inflationary expansion in Planck units, yielding $n \approx 184$.
  \item \textbf{Bekenstein--Hawking consistency.} The same $n$, applied
        to a solar-mass black hole, recovers the Bekenstein--Hawking
        entropy up to the factor $4\pi$ fixed by spherical symmetry.
\end{enumerate}

The saddle-point equation~\eqref{eq:saddle} provides a third,
purely information-theoretic route to the same number. If the solution
of equation~\eqref{eq:saddle} yields $n^* \approx 184$, the three
routes converge on the same value from completely independent
directions. That convergence is not a coincidence. It is the theory
being true.

% ============================================================
\section{Recovering Known Physics}
\label{sec:recovery}

We summarise the derivation chain, each step of which is established in
the corresponding paper of this series.

\subsection{Spacetime Geometry (Papers IV--V)}

With $n$ bits there are $2^n$ distinguishable spatial configurations and
$n \ln n$ bit-flip steps to reach entropy saturation from a zero-entropy
initial state. Their ratio defines the aspect ratio:
\[
\mathcal{A}(n) = \frac{2^n}{n \ln n}.
\]
Elevating from the natural 1D planar geometry to a spherically symmetric
3D topology (the unique minimal-boundary extension) multiplies by $4\pi
R^2$, recovering the observable universe's spatial geometry. The two
extreme solutions of General Relativity --- De Sitter vacuum ($m=0$,
all bits free) and the Schwarzschild singularity ($m=n$, all bits bound
into one structure) --- are the two extreme values of the counting
equation:
\[
R(t) = \frac{(n - m(t)) + k(t)}{n},
\]
where $m(t)$ is the number of bits consumed by matter structures and
$k(t)$ is the count of composite entities. The full Friedmann equation
is recovered under the identifications of Paper~V \cite{meskanen2026v}.

\subsection{Black Hole Thermodynamics (Paper IV)}

Identifying $n_{\mathrm{bh}} = \log_2(r_s / \ell_P)$ for a black hole
of Schwarzschild radius $r_s$ and applying the same spatial resolution
argument yields a surface area that matches Bekenstein--Hawking exactly
once spherical symmetry is imposed. No constants of General Relativity
are introduced; they enter only in the conversion to physical units.

\subsection{Quantum Mechanics and the Born Rule (Paper VI)}

The wave-like character of matter follows from the spectral cost term
$C_\psi$. The prior $\mathbb{P}(\psi) \propto 2^{-C_s(\psi)}$ suppresses
high-frequency (rough) wavefunctions exponentially, exactly as the
Euclidean path integral suppresses high-action histories via
$e^{-S_E}$. Paper~VI \cite{meskanen2026vi} demonstrates that $C_s$ and
$S_E$ are proportional across ensembles of minisuperspace histories
(Spearman $\rho > 0.96$, surviving confound control). The Born rule
emerges as the rasterisation of a compressed description onto a finite
discrete measurement: the probability of an outcome is the fraction of
the compressed bit budget that the outcome occupies, which is $|\psi|^2$
in the spectral representation.

\subsection{Classical Geodesics (Paper VII)}

Paper~VII \cite{meskanen2026vii} shows that minimising $C_s$ in
Wheeler--DeWitt minisuperspace selects the same trajectories as
minimising the Euclidean action. The Hartle--Hawking no-boundary
instanton is the joint global minimum of both measures. Classical
geodesics are not the paths that extremise an action postulated from
outside; they are the paths of shortest description in the observer's
natural basis.

\subsection{Cosmological Expansion (Paper VIII)}

Paper~VIII \cite{meskanen2026viii} applies the lognormal emergence law
of Paper~III \cite{meskanen2003} to the cosmological scale. Without any
hard-coded physics, cosmological constant, or dark energy, the expansion
profile of $\Lambda$CDM is recovered: exponential early expansion,
slowdown, and late-time accelerating expansion. The cosmological
constant is not a parameter; it is the fraction of the bit budget that
remains unbound at the current epoch.

% ============================================================
\section{Discussion}
\label{sec:discussion}

\subsection{Why This Is a Zero-Parameter Theory}

The theory has one structural element: the observer. From the observer
axiom, equation~\eqref{eq:master} follows by Solomonoff induction. The
apparent parameter $\lambda$ is a units choice, not a physical constant.
The apparent parameter $n$ is the saddle point of
equation~\eqref{eq:saddle}, determined self-consistently by the
observer's own existence probability. No constants of nature are
introduced by hand. The constants that appear in physical units
($G$, $c$, $\hbar$, $k_B$) enter only when converting the
dimensionless output of the theory into the unit systems of existing
measurement apparatus.

\subsection{Why We Find Ourselves in a Finite Universe}

Infinite universes are not forbidden. They have infinite description
cost and therefore measure zero under equation~\eqref{eq:master}.
Chaotic universes are not forbidden. They compress poorly and therefore
contribute negligibly to $Z_O$. We find ourselves in a finite, ordered,
geometric universe not because infinity or chaos are impossible, but
because they are astronomically expensive to describe, and we dominate
the measure of compressed descriptions.

\subsection{Why the Universe Waves}

The wave-like character of matter is the signature of spectral
compression. An observer whose description is dominated by the spectral
cost term $C_\psi$ perceives her environment as wave-like because she is
sampling the output of a compression codec. The wavefunction is not a
physical field propagating through spacetime; it is the compressed
representation of the observer's own information structure. We observe
waves for the same reason that pixels in a compressed video wave:
smooth, periodic functions are the cheapest basis for approximating any
sufficiently regular signal.

\subsection{Why There Is No Missing Bridge Between GR and QM}

The standard framing of quantum gravity assumes two separate theories
that must be reconciled. In the present framework there is no gap to
bridge. General Relativity and quantum mechanics are two limiting regimes
of the single cost functional $\mathcal{C}_O$, dominant at different
scales. The ``bridge'' is equation~\eqref{eq:master} itself, read at
the scale where $C_\psi$ and $C_G$ are comparable --- the Planck scale,
where the discrete structure of $C_D[n]$ becomes relevant. There is no
quantum gravity theory to be discovered. There is only one equation,
read at three different scales.

\subsection{The Role of the Observer}

The observer is not an afterthought appended to a physical theory. She
is the only first principle. The universe does not exist independently
and then happen to contain observers. Observers exist, and the universe
is the minimal compressed description that contains them. This is not
solipsism: the same argument that establishes Alice's existence
establishes the existence of every other observer-compatible structure in
the high-probability histories. We are not alone in the measure; we are
the dominant contributors to it. The observable universe is the intersection
of the observers.

% ============================================================
\section{Open Questions}
\label{sec:open}

\subsection{Solving the Saddle-Point Equation}

The central open problem is the analytic or numerical solution of
equation~\eqref{eq:saddle}. If $n^* \approx 184$ emerges from this
equation, the three independent derivations of $n$ converge and the
theory is complete. The solution requires an explicit model of $Z_O(n)$
--- the partition function of observer-compatible histories as a function
of $n$ --- which demands a precise definition of observer complexity in
terms of the spectral cost $C_s$.

\subsection{Geometry as Boundary Condition}

The D-$\psi$-G trinity is forced by the requirement that observers have
well-defined boundaries. But the derivation of the correct topology ---
why the boundary is a 3-sphere rather than some other manifold --- has
not been made fully rigorous. The spherical topology emerges as the
unique minimal-boundary extension of the 1D radial encoding, but a
complete proof that this topology is forced (rather than merely natural)
remains to be written.

\subsection{The Stress-Energy Content}

The present framework recovers the vacuum solutions of General Relativity
and the Friedmann equation under the identification of matter as bound
bit structures. A complete derivation of the Einstein field equations
with a stress-energy tensor, from the counting equation alone, has not
yet been carried out.

% ============================================================
\section{Conclusion}
\label{sec:conclusion}

We have shown that a single equation --- the Solomonoff-weighted measure
over observer-compatible histories --- is sufficient to derive the
structure of spacetime, the origin of quantum mechanics, and the
unification of General Relativity with quantum field theory. The
equation has no free parameters: the apparent parameter $\lambda$ is a
units choice, and the apparent parameter $n$ is the saddle point of the
observer's own existence probability. The value $n^* \approx 184$
predicted by the saddle-point condition is consistent with independent
derivations from inflationary cosmology, black hole thermodynamics, and
the e-fold count of inflation.

The D-$\psi$-G trinity is not assumed. It is forced by the requirement
that any self-describing observer must have boundaries, and that
geometric boundaries are cheaper to describe than spectral ones. We find
ourselves in a geometric, wave-like, finite universe because that is the
most compressible universe consistent with our existence.

The universe is not governed by laws. It is compressed information. The
laws are what compression looks like from the inside.

\printbibliography

\end{document}
